\section{Experiments}\label{sec:experiments}

\subsection{Method}\label{sec:method}

The four sampling choices are interchangeable components of a single filter loop, each selected by name, so a
configuration is a set of names and a particle count. The same loop serves both problems; what differs is the
state a particle carries, a pose in localization and a pose with its own map in SLAM.

Two simulated worlds are used, both taken unchanged from the reference implementations. The \emph{localization
world} is a \SI{10}{\meter} square, cyclic at its edges, with four landmarks. The robot takes 30 steps of
\SI{0.25}{\meter} and \SI{0.02}{\radian}, with motion noise of \SI{0.005}{\meter} and \SI{0.002}{\radian}, and
measures range and bearing to all four landmarks with noise of \SI{0.2}{\meter} and \SI{0.05}{\radian}.
Particles start uniformly over the world and all headings, so the opening steps are global localization; the
learning experiments add a setting in which they start around the true initial pose, so the filter tracks
instead. The \emph{SLAM world} is the FastSLAM scenario of its reference implementation: 500 steps among eight
landmarks, observations within \SI{20}{\meter} with noise of \SI{0.3}{\meter} and \SI{2}{\degree}, and controls
corrupted by noise and a constant yaw-rate offset.

A run is one configuration on one seed, and the seed fixes the world, the trajectory and the filter's randomness.
The same seeds are used for every variant in a comparison, so two variants differ only in the choice being
varied. Every configuration is run on twenty seeds, and every table reports the mean and standard deviation over
them. A difference is called clear when the 95\% confidence intervals of the two means do not overlap, and it is
only taken between variants sharing the value of the other varied setting.

What is measured: accuracy, as position RMSE and as ATE after the best rigid alignment, plus map error in SLAM;
degeneracy, as the mean effective sample size and the fraction of distinct particles surviving a resampling step;
consistency, as the median NEES, about 3 when the particle spread matches the error; cost, as runtime per step;
and fit, as the log likelihood \eqref{eq:evidence} of the observations.

Reusing an implementation means checking it. Four errors were found in the upstream code, among them a FastSLAM
2.0 update that never draws from its proposal, and each would have biased a comparison below. They are listed in
\cref{app:upstream}; all results use the corrected code.

\subsection{Sampling}\label{sec:results-sampling}

Two experiments. Each varies one sampling choice alone; the proposal and the number of particles are held fixed
throughout, at the motion model and at the counts given below.

\subsubsection{Resampling scheme}

\paragraph{Setup.} The four schemes at 50, 200 and 1000 particles in localization, with the motion-model proposal
and resampling at every step so that only the scheme differs (240 runs); and the same four in FastSLAM 1.0 at 10
and 50 particles (160 runs).

\paragraph{Result.} In localization the four are equally accurate: at every particle count their errors differ by
less than the seed-to-seed spread (\cref{fig:E01_resampling_scheme_localization-metrics}). They differ clearly in
how many distinct particles survive a resampling step, in the expected order: with 1000 particles systematic
keeps 72\% distinct, stratified 67\%, residual 65\% and multinomial 54\%, and the order holds at 50 and 200. In
FastSLAM that difference reaches the estimate, because a particle that is lost takes its map with it
(\cref{tab:E01_resampling_scheme_slam}): with 10 particles systematic gives an ATE of \SI{0.23}{\meter} against
\SI{0.32}{\meter} for multinomial, and a map error of \SI{0.105}{\meter} against \SI{0.139}{\meter}. The runtimes
differ by an order of magnitude, but they reflect the reused implementations, three of which cost $O(N^2)$ where
all four can be $O(N)$.

\experimentfigure{E01_resampling_scheme_localization}{metrics}{Resampling schemes in localization, against the number of particles. The four are equally accurate; they differ in how many distinct particles survive.}
\experimenttable{E01_resampling_scheme_slam}{Resampling schemes in FastSLAM 1.0, where each particle carries its own map.}

\subsubsection{When to resample}

\paragraph{Setup.} Eight rules in localization with 100 particles, systematic resampling and the motion-model
proposal: at every step, never, an ESS threshold of 0.2, 0.5 or 0.8, and a largest-weight threshold of 0.1, 0.2
or 0.5 (160 runs).

\paragraph{Result.} Not resampling at all triples the error, \SI{1.39}{\meter} against roughly \SI{0.46}{\meter}
(\cref{fig:E02_when_to_resample-metrics}). Every rule that does resample is equally accurate, whether it acts at
every step or only when the effective sample size or the largest weight crosses a threshold. The thresholds
change how often resampling happens, between a third of the steps and all of them, not how accurate the filter
is; resampling less often leaves each step with more degenerate weights and fewer distinct particles afterwards.

\experimentfigure{E02_when_to_resample}{metrics}{When to resample, in localization with 100 particles. Never resampling stands apart; the rules that do resample are indistinguishable.}

\subsection{Parameter learning}\label{sec:results-learning}

Four experiments. All of them sweep one noise parameter while the other three are held at the values the world
uses, scoring every value by \eqref{eq:evidence}. Unless stated otherwise the filter uses the motion-model
proposal, systematic resampling, an ESS threshold of 0.5 and 2000 particles, and starts around the true initial
pose.

\subsubsection{Recovering the four noise parameters}

\paragraph{Setup.} Each of the four swept in turn: the range noise over nine values (180 runs), the bearing noise
over nine (180 runs), the forward motion noise over seven (140 runs) and the turn noise over seven (140 runs).

\paragraph{Result.} The two measurement parameters come back. Both sweeps have a clear maximum at the true value:
\SI{0.19}{\meter} against a true \SI{0.2}{\meter} (\cref{fig:E11_learn_measurement_range-metrics}) and
\SI{0.047}{\radian} against a true \SI{0.05}{\radian} (\cref{tab:E11_learn_measurement_bearing}), both within one
step of the sweep. Too little assumed noise makes the measurements that arrived look impossible and the score
falls steeply, from 205 at \SI{0.19}{\meter} to $-252$ at \SI{0.08}{\meter}; too much makes every pose explain
them equally well and the score decays slowly, to 98 at \SI{0.8}{\meter}. The penalty is one-sided, as the model
predicts: for a Gaussian residual of true width $\sigma^*$ scored under $\sigma$, the expected score is
$-\log\sigma - \sigma^{*2}/2\sigma^2$ up to a constant, maximized at $\sigma^*$ but falling like $\sigma^{-2}$
below it and only like $\log\sigma$ above. With 120 range residuals that expression predicts the whole measured
curve with nothing fitted: at 0.142, 0.253, 0.45 and \SI{0.8}{\meter} it gives $-18$, $-6$, $-49$ and $-110$
relative to the peak, against measured $-20$, $-5$, $-47$ and $-108$. It parts company only at the far left,
predicting $-205$ where $-457$ is measured, because there the weights degenerate and the estimator's downward
bias takes over.

The two motion parameters do not come back. Both sweeps are flat, the log likelihood holding at 205 to 206 while
the assumed noise varies by a factor of 30 (\cref{tab:E11_learn_motion_forward,tab:E11_learn_motion_turn}). Their
true values, \SI{0.005}{\meter} and \SI{0.002}{\radian}, are small next to the measurement noise the filter
corrects with at every step, so the observations barely depend on them.

Each seed is a separate recording, so an estimate can also be formed from each one alone. All 20 recordings pick
\SI{0.047}{\radian} for the bearing noise and 17 of 20 pick \SI{0.19}{\meter} for the range noise; the 20
estimates of the turn noise spread over the whole sweep. Agreement between recordings is therefore a usable test
of whether a parameter is being estimated at all, and it also tests the tying assumption, since a noise level
that changed between runs would scatter even where the likelihood is sharp.

\experimentfigure{E11_learn_measurement_range}{metrics}{The assumed range noise: the log likelihood peaks at the value the world uses, \SI{0.2}{\meter}, and the error is lowest at the same place.}
\experimenttable{E11_learn_measurement_bearing}{Sweeping the assumed bearing noise; the world uses \SI{0.05}{\radian}.}
\experimenttable{E11_learn_motion_forward}{Sweeping the assumed forward motion noise; the world uses \SI{0.005}{\meter}.}
\experimenttable{E11_learn_motion_turn}{Sweeping the assumed turn noise over 30 steps; the world uses \SI{0.002}{\radian}.}

\subsubsection{How much data identifiability needs}

\paragraph{Setup.} The turn sweep repeated on a 400-step trajectory instead of 30, with nothing else changed,
at 1000 particles (140 runs).

\paragraph{Result.} The flat curve becomes one with a maximum at the true value, and an assumed drift ten times
too small is rejected by thousands of nats (\cref{fig:E11_learn_motion_turn_long-metrics}). The 20 recordings
concentrate again instead of spreading. Identifiability is therefore a property of the data rather than of the
method: because the parameters are tied across time slices, the curvature of the likelihood grows with the number
of steps, as \cref{sec:background} predicts.

\experimentfigure{E11_learn_motion_turn_long}{metrics}{The same turn sweep over 400 steps: enough steps identify a value that 30 steps cannot.}

\subsubsection{Whether the filter has to be tracking}

\paragraph{Setup.} The range sweep run twice, once from a known start pose and once from the uniform start used
by the sampling experiments (360 runs).

\paragraph{Result.} From the uniform start the maximum disappears
(\cref{fig:E11_start_regime-metrics}). With the true, small noise the filter does not find the robot within 30
steps from a uniform prior: the position error is \SI{2.3}{\meter} against \SI{0.063}{\meter} from a known start,
and the log likelihood is four orders of magnitude lower. The score then keeps improving as the assumed noise is
widened, because a wider noise lets the particles spread enough to be useful, and only 2 of 20 recordings agree.
What the likelihood reports in this regime is which setting lets this filter work, not which setting generated
the data.

\experimentfigure{E11_start_regime}{metrics}{The range sweep from a known start and from the uniform start of the sampling experiments. Note the scale: from the uniform start the estimate is four orders of magnitude lower and has no maximum.}

\subsubsection{What the estimate needs from the inference}

\paragraph{Setup.} The range sweep repeated at 25, 100, 500 and 2000 particles (720 runs), and again under each
of the four resampling schemes at 200 particles (720 runs).

\paragraph{Result.} At the true value the estimated log likelihood rises from $28 \pm 316$ with 25 particles to
$205 \pm 13$ with 2000, and agreement between recordings rises with it, from none at 25 particles to 17 of 20 at
2000 (\cref{fig:E11_estimator_particles-best_per_seed}): both the bias and the spread shrink as the approximation
improves. The resampling scheme makes no difference: at 200 particles all four place the maximum at
\SI{0.19}{\meter}, with 14 or 15 of 20 recordings agreeing (\cref{tab:best-E11_estimator_resampler}).

\experimentfigure{E11_estimator_particles}{best_per_seed}{The range noise each recording picks out on its own, against the number of particles.}
\experimentbest{E11_estimator_resampler}{The range noise picked out under each resampling scheme, with 200 particles.}
